% Korean edition of sections/s2_0_parent.tex.
\seclabel
\claim{\secnum{2}{sec:layers}\quad 실서버에서 측정한 프로토콜 5개 계층이 래퍼의 매개변수를 정한다}

\lead{래퍼가 쓰는 모든 타임아웃, 간격, 크기는 측정값이나 명시된 규칙에 연결된다.} 측정은 2026년 11월 24일 05:38\textasciitilde 10:34 KST에 실서버를 대상으로 했다. 예외는 2026년 10월 28일 브라우저 기록에서 얻은 3\,s 폴링 간격과, 주최측 이메일[5]에서 얻은 하루 주문 한도다. 한도를 측정했거나 명시받은 곳에서는 래퍼가 그 한도 안쪽으로 여유를 두고 자기 매개변수를 정한다(\ref{sec:design}장). 주최측 문서가 없으니 측정하지 않으면 추측해야 하는데, \ref{sec:queries}절에서 보듯 서버는 잘못된 추측에도 HTTP 200으로 답한다. 각 절은 계층 하나를 측정하고 그 계층이 정하는 매개변수로 끝난다. \ref{sec:transport}는 전송(재사용 연결, 106.5\ms{} 대비 16.0\ms), \ref{sec:auth}는 인증(900\,s 액세스 토큰과 7일 쿠키 묶음), \ref{sec:queries}는 조회(날짜는 서버에서 받음), \ref{sec:feed}는 실시간 피드(3\,s 폴링과 델타), \ref{sec:orders}는 주문(넣기, 확인, 크기)이다. 서버는 이 동작을 언제든 바꿀 수 있다. \ref{sec:validity}절은 래퍼가 주문 전마다 검사하는 항목을 나열한다.
